\section{1. Introduction}

Heterogeneous UAV teams combine sensing, communication, and task-execution roles. Their cooperative capability depends not only on each vehicle's kinematics but also on whether task-relevant information can traverse legal sensing and communication relations. Graph-based multi-agent reinforcement learning (MARL) provides representations for structured interaction [2,3,17--21], and graph neural policies have been applied to multi-UAV search, tracking, and swarm control [6,7]. A relay-node failure, however, changes the relation structure itself. When a failed relay removes a multi-hop route while a direct Scout--Attacker route remains legal, the problem is not complete information loss. It is topology/path and task-support reconfiguration under a strict decentralized information boundary.

This setting exposes a training-distribution question that is separate from graph encoding. Robust MARL studies partner, model, and adversarial variation [4], while distributionally robust reinforcement learning studies explicitly specified distributional ambiguity [5]. Group-robust learning, active domain randomization, curriculum learning, level replay, and unsupervised environment design instead study how training exposure can be allocated across groups, parameters, or generated tasks [11--13,22--26]. These approaches motivate adaptive exposure but also illustrate that changing a training distribution can change the optimization problem. They do not make a bounded heuristic sampler a minimax solver or provide a general distributionally robust guarantee for the frozen relay-failure setting.

We therefore study a deliberately narrow comparison: with the Single-Graph MAPPO backbone, PPO objective, reward, topology-perturbation support, nominal-condition exposure, budget, and paired evaluation protocol held fixed, does bounded adaptive reweighting of six predefined relay-failure groups improve mission performance relative to uniform reweighting? We ask both whether gains appear on canonical F0 and across the frozen perturbation family, and whether their direction persists in a separately completed training cohort. The second question is part of the empirical result, not a post hoc exception.

To answer this question, we introduce bounded adaptive topology-perturbation reweighting Single-Graph MAPPO (DRTP-SG-MAPPO; DRTP). DRTP retains a 50% nominal-condition anchor and redistributes the remaining failure mass across six groups using return gaps relative to nominal competence, smoothing, and bounded-simplex projection. Its primary comparator, Uniform Topology Randomization Single-Graph MAPPO (UTR-SG-MAPPO; UTR), has the same 116,728-parameter actor--critic, PPO settings, environment, reward, perturbation support, budget, and execution-time information boundary. The comparison identifies an empirical difference between bounded adaptive and uniform reweighting in this task; it does not establish necessity relative to every fixed non-uniform distribution.

Our contributions are threefold.

1. We formulate a controlled three-UAV relay-failure coordination task in which a failure can remove Relay-mediated sensing and communication relations while preserving a legal direct Scout--Attacker path. The task makes topology/path and task-support reconfiguration observable without injecting failure labels or global topology truth into decentralized actors.
2. We define DRTP as a bounded training-time reweighting controller over six frozen failure/topology groups. It leaves the actor--critic architecture, PPO objective, reward, and execution-time information unchanged, concentrating the method contrast on the training distribution.
3. We report a reliability-aware evaluation. A prospective, parameter- and exposure-matched five-seed 10M formal cohort favored DRTP on F0, \texttt{J_pert,mean\texttt{, and \texttt{J_pert,worst\texttt{, with a timeout--collision trade-off. A separately completed three-method cohort reversed the DRTP--UTR direction on both evaluation tapes. We report both strata rather than claim stable superiority or strict OOD generalization.

\section{3. Related work and positioning}

\subsection{3.1 Graph-structured multi-agent learning}

Communication-learning methods include continuous communication, differentiable inter-agent learning, targeted message routing, and attentional communication [17--20]. Attention can also be used in a centralized critic to select information relevant to each decentralized agent [21]. Topology-aware policy-gradient work and UAV graph-MARL studies further show that agent relations can be represented explicitly in cooperation and swarm-control problems [3,6,7]. These studies motivate the fixed Single-Graph encoder and legal information representation in our task. DRTP does not claim a new communication architecture; it holds the actor--critic fixed and studies how topology perturbations are distributed during training.

\subsection{3.2 Robust learning and adaptive training distributions}

Robust MARL includes minimax formulations for partner or environmental variation [4], and distributionally robust Q-learning formalizes ambiguity in a transition or environment distribution [5]. Group-robust reweighting [11], active domain randomization [12], curricula [13,25], prioritized level replay [22], and unsupervised environment design [23,24] demonstrate several ways to alter training exposure. Robust on-policy sampling likewise distinguishes the data distribution from the trajectory-collection schedule [26]. DRTP shares only the empirical premise that uniform exposure can underrepresent conditions that are difficult for a current policy. It is neither a minimax solver nor a certified group-DRO method, and it makes no worst-case guarantee. Its distinction is a reproducible, bounded six-group sampler evaluated against matched uniform sampling in one fixed relay-failure contract.

\subsection{3.3 Communication-limited UAV cooperation and fault-related topology changes}

MARL has been applied to UAV swarm communication under jamming [8,27,28], multi-hop relay communication [9], network recovery after UAV failure [29], and cooperative trajectory or resource allocation [30]. Communication-network and consensus studies also describe how link or node failures can require path or topology reconfiguration [15,16]. These studies establish application relevance, but their action spaces, communication objectives, reward definitions, learning algorithms, and information boundaries differ from the present relay-failure coordination task. They are used for problem positioning, not converted into nominally fair performance baselines.

\subsection{3.4 Comparator boundary}

The paper's primary causal comparison is UTR versus DRTP under a common SG-MAPPO backbone. MAPPO-NoGraph is reported as an external performance reference because it shares the task, seeds, budget, and evaluation tape but differs in graph-message input and parameter count. The independent cohort additionally includes a fixed non-uniform sampler, yet it has a separate cohort and evaluation tape and does not reproduce the formal DRTP direction. Neither comparison can be backfilled as evidence that online adaptation is necessary relative to every static sampler. This evidence hierarchy separates the controlled within-contract claim from external performance context and from the reliability limitation.